\section{Abelian categories}

We are now prepared to begin defining \emph{abelian categories}. We already know of one requirement they must satisfy: an abelian category must be an additive category. This section will look for other properties they should satisfy in order for us to have a good theory of homological algebra.

The first requirement we must make is the existence of kernels. Kernels are usually constructed as \emph{subsets} of $R$-modules or groups, so to generalize their definition to arbitrary categories, we must explore different characterizations of kernels. The most natural characterization is by way of a universal property.

To see how kernels arise by a universal property, let $\varphi : M \to N$ be $R$-linear. Then the kernel $\Ker \varphi $ is a \emph{submodule} of $M$. This means that there is an inclusion map $i:\Ker \varphi \hookrightarrow M$ which is also $R$-linear. This inclusion map satsifies an interesting property. Every $m\in \Ker \varphi $ gets mapped to 0 by $\varphi $ by the definition of a kernel, meaning $\varphi \circ i=0$, since the inclusion simply restricts $\varphi $ to the kernel. Additionally, the kernel is \emph{universal} with respect to this property, meaning for any $\psi : P \to M$ such that $\varphi \circ \psi =0$, the map $\psi $ factors uniquely through the kernel of $\varphi $. Alternatively, given any diagram of solid arrows
\[\begin{tikzcd}[row sep = 2.5em, column sep = 2.5em]
P \ar[" \psi  ", r] \ar[" 0 ", bend left, rr] \ar[dashed, dr]  & M \ar[" \varphi  ", r]  & N \\
 & \Ker \varphi \ar[" i "', u]  & 
\end{tikzcd}\]
there is a unique dashed arrow as indicated, rendering the diagram commutative.

Since zero morphisms exist in $\Ab $-categories, we have an immediate generalization of the notion of kernel to $\Ab $-categories.

\begin{definition}\label{4.1}
	Let $\mathcal{C} $ be an $\Ab $-category, and $\varphi : x \to y$ a morphism in $\mathcal{C} $. A kernel $\Ker \varphi $, if it exists, is an object  of $\mathcal{C} $, together with a morphism $\ker \varphi : \Ker \varphi  \to x$, such that for any diagram of solid arrows
	\[\begin{tikzcd}[row sep = 2.5em, column sep = 2.5em]
	z \ar[r] \ar[" 0 ", bend left, rr] \ar[dashed, dr]  & x \ar[" \varphi  ", r]  & y \\
	 & \Ker \varphi \ar[" \ker \varphi  "', u]  & 
	\end{tikzcd}\]
	there is a unique dashed arrow as indicated, rendering the diagram commutative.
\end{definition}

Note the interesting shift in notation. The morphism $\Ker \varphi \to x$ is now denoted by $\ker \varphi $. Since most of the structure in arbitrary categories is carried by the morphisms, when working outside of $\Mod_R$, it usually makes sense to refer to the \emph{inclusion morphism} as the kernel of $\varphi $, hence our notation. We use the capital K in $\Ker \varphi $ to denote the object, and the lowercase k in $\ker \varphi $ to denote the morphism.

\begin{remark}\label{4.2}
Kernels also have a dual, called the \emph{cokernel}. Given some $\varphi : x \to y$ in an $\Ab$-category $\mathcal{C} $, a \emph{cokernel} $\Coker \varphi $ of $\varphi $ is an object, together with a morphism $\coker \varphi : y \to \Coker\varphi $ with $\coker \varphi \circ \varphi =0$, such that for any diagram of solid arrows
\[\begin{tikzcd}[row sep = 2.5em, column sep = 2.5em]
 & \Coker \varphi \ar[dashed, dr]  &  \\
x \ar[" \varphi  "', r] \ar[" 0 "', bend right, rr]  & y \ar[r] \ar[" \coker \varphi  ", u]  & z
\end{tikzcd}\]
there is a unique dashed arrow as indicated, rendering the diagram commutative.
\end{remark}


We will also need a good notion of an \emph{image} object $\Im \varphi $ for any morphism $\varphi : x \to y$ in an additive category, and we will also need a way to construct the \emph{quotient} of the image and the kernel, such that we can obtain homology. Let's begin by considering quotients.

Let $\varphi : M \to N$ be $R$-linear. The cokernel $\Coker \varphi $ can be constructed as the quotient module
\[
	\Coker\varphi \cong \frac{N}{\Im \varphi } 
.\]
This follows since this object is initial with respect to mapping $\Im \varphi $ to 0. Note that since this is constructed as a quotient module, we might want to think of cokernels as quotients, giving us an abstraction of cokernels. Since $\Im d_n \subseteq \Ker d_{n-1} $ in a chain complex $C_{\bullet} $, we can thus view the homology $H_n(C_{\bullet} )$ as the cokernel of $d_n$, viewed as a map $C_{n} \to \Ker d_{n-1} $, rather than a map $C_n\to C_{n-1} $.

This gives a purely diagram-theoretic construction of homology. Let $C_{\bullet} $ be a chain complex, and since $d_{n-1} d_n=0$, the map $d_n$ factors through the kernel as indicated by the left triangle
\[\begin{tikzcd}[row sep = 2.5em, column sep = 2.5em]
C_n \ar[" d_n ", r] \ar[" d_n' "', dashed, dr]  & C_{n-1} \ar[" d_{n-1}  ", r]  & C_{n-2}  \\
 & \Ker d_{n-1} \ar[" \ker d_{n-1}  "', u] \ar[" \coker d_n'  "', r]  & \Coker d_n' \ar[dashed, u] 
\end{tikzcd}\]
By commutativity of the left triangle, $d_{n-1} \circ \ker d_n \circ d'_n=0$, and thus the second two maps factor through the cokernel. By our knowledge of kernels in $\Mod_R$, we know that $\Im d_n'=\Im d_n$, when viewed as submodules of $C_{n-1} $. Thus, if we demand that kernels and cokernels exist, we are guarenteed a purely categorical construction of homology by taking $H_n(C_{\bullet} ) \coloneqq \Coker d_n'$.

There is one final requirement of abelian categories, which is motivated by a particular property of $\Mod_R$ that characterizes monics and epis. A monic map $M\to N$ is simply an injective map, meaning it identifies $M$ with a subspace $M \cong \Im \varphi \subseteq N$. Since all subspaces are kernels, $\Im \varphi =\Ker \psi $ for some $\psi $, meaning $M$ can be identified with the kernel of $\psi $, and thus $\varphi $ \emph{is a kernel}. Thus, monics are kernels. Conversely, epis coencide with cokernels. This property makes kernels and cokernels extremely well-behaved, and will help us develop a lot of important theory, such as a categorical definition of images and coimages. We thus demand that abelian categories have this property as well.

We have now identified all the properties an abelian category must have, in order to have a well-behaved theory of homology. We summarize with

\begin{definition}\label{4.3}
	Let $\mathcal{C} $ be an $\Ab $-category. Then $\mathcal{C} $ is an \emph{abelian category} if
	\begin{enumerate}
		\item $\mathcal{C} $ is additive (\cref{3.4});
		\item every morphism in $\mathcal{C} $ has a kernel and a cokernel;
		\item every monic arrow is a kernel, and every epi arrow is a cokernel.
	\end{enumerate}
\end{definition}
